\chapter[Installation and Demonstration]{Installation, Execution, and Demonstration Guide}
\label{chap:installation-demo}

\section{System requirements}

The implementation was developed and tested in a Linux/WSL-style environment. The exact
gas measurements reported in Chapter~\ref{chap:experimental-results} were obtained on a
local Hardhat network. Reproducing the project therefore requires both the web stack and
the native cryptographic tooling.

The required tools are:

\begin{itemize}
  \item Node.js, version 22 or later;
  \item npm, used by the Next.js frontend and the Hardhat project;
  \item Rust and Cargo, used to build the native precontract and dispute CLIs;
  \item \texttt{wasm-pack}, needed only when rebuilding the frontend WebAssembly bundle;
  \item \texttt{sqlite3}, used to initialize the local database;
  \item \texttt{pnpm} and Foundry tooling, used by the Alto ERC-4337 bundler setup;
  \item a Unix-compatible shell environment, preferably Linux or WSL.
\end{itemize}

The project contains an installation script that checks most of these dependencies and
installs the JavaScript packages. The script is useful, but the manual commands below are
also given so that the deployment can be reproduced and debugged step by step.

\section{Repository setup}

The repository should be cloned with its submodule, because the Alto bundler is stored in
\path{bundler-alto} as a Git submodule.

\begin{verbatim}
git clone --recurse-submodules <repository-url>
cd sox_implementation-D0AB
git submodule update --init --recursive
\end{verbatim}

From the repository root, the standard installation command is:

\begin{verbatim}
chmod +x install.sh
./install.sh
\end{verbatim}

The script installs the root Next.js dependencies, the Electron dependencies, the Hardhat
dependencies, initializes the SQLite database if needed, and sets up the Alto bundler. It
also builds the native \texttt{precontract\_cli} if Cargo is available. If the script is
not used, the minimum manual installation is:

\begin{verbatim}
npm install
cd src/hardhat && npm install && cd ../..
cd desktop && npm install && cd ..
\end{verbatim}

The local database is stored in \path{src/app/db/sox.sqlite}. To create a fresh database,
remove the old file and replay the schema:

\begin{verbatim}
rm -f src/app/db/sox.sqlite
sqlite3 src/app/db/sox.sqlite < src/app/db/init.sql
\end{verbatim}

This resets the local application state. It does not reset the blockchain node; restarting
Hardhat is the clean way to reset the chain.

\section{Building the Rust and WebAssembly components}

The project uses Rust in two different ways. The frontend imports a WebAssembly bundle from
\path{src/app/lib/crypto_lib}. The large-file benchmarks use native Rust binaries from
\path{src/wasm/target/release}. These are separate build targets and should not be
confused.

To rebuild all native CLI binaries used by the benchmarks, run:

\begin{verbatim}
cd src/wasm
cargo build --release --bins
cd ../..
\end{verbatim}

The most important binaries are:

\begin{itemize}
  \item \texttt{precontract\_cli}, which generates precontract artifacts;
  \item \texttt{check\_precontract\_cli}, which verifies received precontract data;
  \item \texttt{compute\_proofs\_cli}, which supports the current proof-generation API;
  \item \texttt{hardcoded\_sha256\_dispute\_cli}, which generates hardcoded dispute data
  for the large-file benchmark.
\end{itemize}

To rebuild the WebAssembly bundle used by the web application, install \texttt{wasm-pack}
if necessary and run:

\begin{verbatim}
cd src/wasm
./deploy.sh
cd ../..
\end{verbatim}

This writes the generated JavaScript and WASM files to
\path{src/app/lib/crypto_lib}. The Hardhat performance tests that import the WASM bundle
expect these files to exist.

\section{Deploying the local smart contract stack}

The local blockchain node must be running before contracts are deployed. In a first
terminal, start Hardhat:

\begin{verbatim}
cd src/hardhat
npx hardhat node
\end{verbatim}

In a second terminal, from the repository root, deploy the complete local contract stack:

\begin{verbatim}
./deploy-all.sh
\end{verbatim}

The deployment script performs three main tasks. It deploys the local EntryPoint, deploys
the simulation contracts needed by Alto, and deploys the SOX libraries, dispute deployers,
clone implementations, and factory. The generated addresses are written to
\path{deployed-contracts.json} and \path{src/deployed-contracts.json}. These files are used
by the web application and the deployment helper.

If the Hardhat node is restarted, the deployed addresses are no longer valid. In that case,
run \texttt{./deploy-all.sh} again before using the application or the bundler.

\section{Running Hardhat, Alto, and the frontend}

The recommended development setup uses separate terminals. Keeping the processes separated
makes errors easier to identify.

\begin{table}[htbp]
  \centering
  \begin{tabularx}{\linewidth}{>{\raggedright\arraybackslash}p{0.18\linewidth}>{\raggedright\arraybackslash}X}
    \toprule
    Terminal & Command \\
    \midrule
    1 & \texttt{cd src/hardhat \&\& npx hardhat node} \\
    2 & \texttt{./deploy-all.sh} after Terminal 1 is ready \\
    3 & \texttt{./run-alto.sh} after deployment succeeds \\
    4 & \texttt{npm run dev} from the repository root \\
    5, optional & \texttt{cd desktop \&\& npm start} after Next.js is running \\
    \bottomrule
  \end{tabularx}
  \caption{Recommended local process layout.}
  \label{tab:local-process-layout}
\end{table}

The web application is then available at \url{http://localhost:3000}. The local Hardhat RPC
is available at \url{http://localhost:8545}. The Alto bundler RPC is available at
\url{http://localhost:4337/rpc}.

The helper script \path{START_ALL.sh} can be used to coordinate the services, but the
manual terminal layout in Table~\ref{tab:local-process-layout} is the most reproducible
method for debugging and for demonstrations.

\section{Running the optimistic demo}

The simplest end-to-end demonstration is an honest optimistic exchange with the normal
sponsored path. A 4 MiB test file can be created with:

\begin{verbatim}
mkdir -p demo_files
truncate -s 4M demo_files/demo_4MiB.bin
\end{verbatim}

With Hardhat, the deployed contracts, Alto, and Next.js running, open the web interface at
\url{http://localhost:3000}. The demo flow is:

\begin{enumerate}
  \item Select the vendor role and create a new precontract.
  \item Choose the buyer address, upload \path{demo_files/demo_4MiB.bin}, and keep the
  normal sponsored mode for the standard demo.
  \item Submit the precontract. The backend calls the native precontract generator and
  stores the ciphertext and metadata.
  \item Switch to the buyer role and open the non-accepted precontract.
  \item Click ``Verify commitment'' to call \path{/api/precontracts/verify}. This checks
  the received ciphertext and commitment data with the native precontract checker.
  \item Accept the precontract. The application marks it as accepted and downloads the
  ciphertext for the buyer.
  \item Switch to the sponsor role and sponsor/deploy the optimistic contract.
  \item Switch to the buyer role and send the buyer payment.
  \item Switch to the vendor role and send the key.
  \item The optimistic contract reaches the complete state if the exchange is honest.
\end{enumerate}

For demonstrations focused on the new optimistic variants, the vendor can select
\simsox{} (\texttt{no\_S\_deposit}), \(S=B\), or \(S=V\) in the precontract creation modal.
The user interface enforces the obvious role constraints, and the corresponding on-chain
path repeats the relevant checks.

\section{Running gas benchmarks}

The Hardhat test suite contains both correctness tests and performance tests. The following
commands are the main tests used for the measurements in this report:

\begin{verbatim}
cd src/hardhat
npx hardhat test test/Phase3Variants.test.ts
npx hardhat test test/Phase3ExhaustiveAndGas.test.ts
npx hardhat test \
  test/performance/HanaComparisonCurrent.test.ts
npx hardhat test \
  test/performance/SpecializedContractArchitecture.test.ts
npx hardhat test \
  test/performance/SpecializedFinalSameScope.test.ts
\end{verbatim}

The first two commands test the implemented variants and the monolithic variant matrix. The
\texttt{HanaComparisonCurrent} benchmark remeasures the same scopes as the initial
implementation. The specialized architecture benchmarks measure the factory, clone,
self-sponsored, and hardcoded paths used in Chapter~\ref{chap:experimental-results}.

Gas reporter output can also be enabled through the environment variable already supported
by the Hardhat configuration:

\begin{verbatim}
REPORT_GAS=true npx hardhat test
\end{verbatim}

For stable comparisons, the repository should be measured on a fresh Hardhat state and with
the same Solidity optimizer settings as the submitted implementation. The Hardhat
configuration uses Solidity 0.8.28, optimizer enabled, \texttt{viaIR}, and optimizer
\texttt{runs = 1}, because the project also targets bytecode-size reduction.

\section{Running large-file benchmarks}

The 1 GiB hardcoded benchmark requires the native Rust binaries. Build them first:

\begin{verbatim}
cd src/wasm
cargo build --release --bin precontract_cli \
  --bin hardcoded_sha256_dispute_cli
cd ../hardhat
\end{verbatim}

Then run the benchmark:

\begin{verbatim}
npx hardhat test \
  test/performance/NativeHardcodedSHA256FullDispute1GB.test.ts
\end{verbatim}

By default, this test uses a 1 GiB file. It creates the file in a temporary directory using
file truncation, runs \texttt{precontract\_cli}, runs the hardcoded dispute CLI, executes
all challenge rounds, performs Step 8, and finalizes the dispute. The temporary directory
is removed at the end of the test.

For a smoke test with the same code path but a smaller file, set the benchmark
size environment variable. For example:

\begin{verbatim}
SOX_FULL_DISPUTE_SIZE_BYTES=16384 npx hardhat test \
  test/performance/NativeHardcodedSHA256FullDispute1GB.test.ts
\end{verbatim}

The 1 GiB run requires enough disk space for the temporary plaintext, ciphertext, and
intermediate artifacts, and it takes significantly longer than the small-file tests. It is
therefore useful to run the small smoke test first when checking a new machine.

\section{Troubleshooting}

Most setup problems come from stale local state. A clean reset usually means restarting the
Hardhat node, redeploying the contracts, restarting Alto, and restarting Next.js.

\begin{table}[htbp]
  \centering
  \small
  \begin{tabularx}{\linewidth}{>{\raggedright\arraybackslash}p{0.32\linewidth}>{\raggedright\arraybackslash}X}
    \toprule
    Symptom & Fix \\
    \midrule
    \texttt{spawn precontract\_cli ENOENT} &
    Build the native binary from \path{src/wasm} with
    \texttt{cargo build --release --bin precontract\_cli}. \\
    Missing \path{deployed-contracts.json} or missing library addresses &
    Start Hardhat and rerun \texttt{./deploy-all.sh}. \\
    Alto says the EntryPoint does not exist &
    The Hardhat chain was probably restarted. Redeploy with \texttt{./deploy-all.sh} and then restart Alto. \\
    Frontend reports \texttt{Failed to fetch} for UserOperations &
    Check that \texttt{./run-alto.sh} is running and that the bundler RPC is reachable at
    \url{http://localhost:4337/rpc}. \\
    \texttt{EACCES} in \path{src/hardhat/artifacts} &
    Some files were created by another user. Restore ownership of
    \path{src/hardhat/artifacts} and \path{src/hardhat/cache}. \\
    Native 1 GiB benchmark skips itself &
    Build both required native binaries: \texttt{precontract\_cli} and
    \texttt{hardcoded\_sha256\_dispute\_cli}. \\
    Browser/WASM struggles on very large files &
    Use the native CLI benchmark path. The 1 GiB result in this report is a native CLI result, not a browser result. \\
    \bottomrule
  \end{tabularx}
  \caption{Common setup and execution issues.}
  \label{tab:troubleshooting}
\end{table}

One additional Alto-specific issue can occur on a fresh local chain: the deterministic
deployment sender used by the bundler may have no balance. On Hardhat, the account can be
funded explicitly through JSON-RPC if needed:

\begin{verbatim}
cat >/tmp/alto-fund.json <<'JSON'
{"jsonrpc":"2.0","id":1,"method":"hardhat_setBalance",
"params":["0x3fAB184622Dc19b6109349B94811493BF2a45362",
"0x8AC7230489E80000"]}
JSON

curl -X POST http://127.0.0.1:8545 \
  -H "Content-Type: application/json" \
  --data @/tmp/alto-fund.json
\end{verbatim}

After this, restart \path{run-alto.sh}. This command is only a local development fix; it is
not part of the SOX protocol.
